\documentclass[11pt]{scrartcl}

\usepackage[stixtwo]{fontsetup}
\usepackage{amsmath}

\begin{document}
\begin{itemize}
	\item \[
		      c_n
		      = \frac{\left(\sqrt{n + \sqrt{n}} - \sqrt{n}\right)\left(\sqrt{n + \sqrt{n}} + \sqrt{n}\right)}{\sqrt{n + \sqrt{n}} + \sqrt{n}}
		      = \frac{\sqrt{n}}{\sqrt{n + \sqrt{n}} + \sqrt{n}}
		      = \frac{1}{\sqrt{1 + \frac{1}{\sqrt{n}}} + 1}
		      \to \frac{1}{2}
	      \]

	\item \[
		      d_n
		      = \frac{-2 \cdot (-32)^n - \frac{1}{3} (-27)^n}{16^n}
		      = (-1)^{n + 1} \left( 2 \cdot 2^n + \frac{1}{3} \left(\frac{27}{16}\right)^n \right)
	      \]
	      A zárójelben álló kifejezés $\infty$-hez tart, az előjel pedig
	      váltakozik: páros $n$-re $d_n \to -\infty$, páratlan $n$-re
	      $d_n \to \infty$. A sorozatnak tehát két torlódási pontja van
	      ($\pm \infty$), így nincs határértéke.
\end{itemize}
\end{document}
